\section{Complete checkers and verification}
\label{sec:implementation}

\subsection{Public contract}

Each compared module reads one JSON object with exactly two fields,
\code{goal} and \code{proof}. The goal must be a closed Eq/$D$ formula.
The proof is a finite tagged tree. The process prints \code{ACCEPT} and
returns success exactly when its inference routine returns that goal
from the empty context, subject to ordinary finite machine resources.
Otherwise it rejects. Unknown tags, extra fields, malformed trees,
noninteger indices, scope violations, and escaped eigenparameters are
rejected.

The mathematical calculus also has arbitrary retained clauses and open
sequents. The public executable interface checks closed theorems, as
does the compared incoming clause checker. Arbitrary open consequences
can be expressed using the derived implication and universal macros;
this does not turn the public JSON contract into an open-sequent API.

\subsection{A compact but explicit binder representation}

The same numerical convention is used in the new checker and both
matched controls. Zero is the integer $0$. A fresh free parameter is a
negative integer level. Positive integers are local bound indices,
counting outwards from $1$. At free scope $k$ and local binder depth $d$,
a scalar integer leaf must lie in the interval $[-k,d]$.

In the raw language, successor is \code{[0,t]}, iteration is
\code{[n,a,h]}, equality is \code{[s,t]}, and $D$ is \code{[B]}.
The expected sort distinguishes a successor from an equation.
Iteration binds one variable in its third child, and $D$ binds two
variables in its child. For $D_{x,y}B$, $y$ is index $1$ and $x$ is
index $2$. This gives one structurally recursive binder-opening function.

Opening replaces the relevant bound index and decreases larger indices
by one. All inserted values are closed with respect to free positive
indices: their free variables are negative levels. Therefore insertion
does not require shifting those values under binders. Their own internal
bound indices stay bound to their own constructors. Retained contexts
also do not need shifting when a fresh negative level is introduced.

This convention is an implementation choice shared by the controls; it
is not credited as a new mathematical simplification. The canonical
control additionally erases the body of an iterator whose rebuilt
counter is literal zero. The new checker performs no such normalization.

\subsection{The inference invariant}

For every successful call to \code{derive(p,g,k)}, every retained clause
is well scoped at $k$, its antecedent is a formula, its consequent is a
one-variable formula template, and the returned value is a well-scoped
formula at $k$. Induction on a successful finite inference trace proves
this invariant.

Reflexivity validates its term. Clause use validates the index and
witness, proves the exact antecedent, and opens the consequent. Eval
destructures two successfully derived formulas as equations. A compound
formula has one child, so equation destructuring rejects it. A term being
unfolded must destructure as a three-child iterator. A nonzero counter
must have the successor shape; the validated raw grammar makes its
two-child test exact. Hence the shorter structural destructuring does
not bypass sort checking.

Paired introduction validates its entire target and witness before
opening the body. Its new hypothesis is retained as a constant clause
with antecedent $0=0$. Paired elimination validates its target before
opening the existential witness. In both branches the fresh negative
level is outside the scope in which the target was validated, preventing
eigenvariable escape. The exact branch result is checked before the
parent result is returned.

The canonical and raw computation/transport controls use the same
contract and invariant. The comparisons count the full validation and
input boundary, not only the small inference kernel.

\subsection{Executed verification and its limits}

The final execution accepted 32 representative positive proofs and verified both raw proof-translation directions with identical end formulas. It rejected 24 adversarial certificates. All 112 public command-line runs, including runs with Python optimization enabled, returned the expected status and output. The canonical control also passed its representative reflexivity, introduction, and canonical-syntax checks.

The accepted corpus includes both iterator computation equations, nested term templates, a used induction hypothesis, clause specialization, compound binder nesting, and capture-avoiding witness insertion. The raw six-to-five translation satisfied the stated bound of at most nine times as many inference nodes on every executed example; this does not bound term-label sizes or JSON length.


The package supplies bidirectional proof transformations for the matched
raw equality interfaces. The mathematical derivations justify them for
all finite derivations; the recorded executable checks exercise concrete
instances, including nesting and parameter binding. Negative examples
target violations that could change the theory: unrestricted compound
induction, wrong Eval sources, escaped fresh parameters, wrong clause
antecedents, and malformed or mistyped public input.

The finite tests establish the stated execution results. They cannot prove
that an arbitrary implementation is sound on every possible certificate.
The universal equivalence is the written proof in this report; no new
proof-assistant verification of these Python modules is asserted.

\subsection{Reproducing the archive}

Use Python 3.12 to run the verification and metric scripts documented in
\code{README.md}. Check a supplied JSON certificate by redirecting it to
\code{python3 code/clausal\_eval.py}. Run \code{latexmk -pdf report.tex}
from the archive root to rebuild this document. The report source, its
inputs, standalone modules, normalized forms, gzip-9 files, test evidence,
and original-source attribution are all included.
